\documentclass[11pt,a4paper]{report}

\usepackage[utf8]{inputenc}
\usepackage[T1]{fontenc}
\usepackage[spanish,es-tabla]{babel}
\usepackage{graphicx}
\usepackage{booktabs}
\usepackage{tabularx}
\usepackage{float}
\usepackage{geometry}
\usepackage{listings}
\usepackage[hidelinks]{hyperref}

\geometry{margin=2.5cm}
\graphicspath{{../lab-evidencias/}}

\lstset{
  language=bash,
  basicstyle=\ttfamily\footnotesize,
  breaklines=true,
  breakatwhitespace=false,
  columns=fullflexible,
  keepspaces=true,
  showstringspaces=false,
  frame=single,
  framerule=0.3pt
}

\begin{document}

\sloppy

\begin{titlepage}
  \centering
  \vspace*{1cm}


  {\large Laboratorio No. 01 --- Base Platform \par}
  \vspace{1.5cm}

  {\large Por: \par}
  {\large Camilo Aguirre \par}
  {\large Juan David Rangel \par}
  \vspace{1cm}

  {\large Profesor: John Alexander Pachón Pinzón \par}
  \vspace{1.5cm}

  {\large Escuela Colombiana de Ingeniería Julio Garavito \par}
  {\large Universidad \par}
  \vspace{0.5cm}

  {\large ISIS AYSR - 1 \par}
  {\large Arquitectura y Servicios de Red \par}
  \vspace{1.5cm}

  {\large Bogotá, D.C., viernes 7 de agosto de 2026 \par}

  \vfill
\end{titlepage}

\tableofcontents
\newpage

\section{Introducción}

Objetivo del laboratorio: instalar y configurar distintas distribuciones de Unix y Windows como parte del montaje de una plataforma, y familiarizarse con el software de virtualización.

\begin{itemize}
  \item \textbf{Contexto:} una empresa opera servicios de TI (estaciones de trabajo cableadas e inalámbricas, servidores físicos y virtualizados, switches L2/L3, routers y nube).
  \item \textbf{Alcance:} instalación de Slackware, Solaris, Windows Server (2 fases), Android y CentOS (solo grupos de 3; este grupo es de 2, no aplica).
  \item \textbf{Herramientas:} VMware, máquinas del laboratorio, ISO descargadas.
\end{itemize}

\section{Marco teórico}

\subsection{Virtualización}

\textbf{¿Qué es un hipervisor?}

Un hipervisor (o monitor de máquina virtual, VMM) es una capa de software que crea y administra máquinas virtuales sobre un mismo hardware físico. Actúa como intermediario entre el hardware y los sistemas operativos invitados: asigna recursos (CPU, memoria, almacenamiento y red), atiende las instrucciones privilegiadas de cada invitado y garantiza el aislamiento entre máquinas: una VM no puede acceder a la memoria ni a los datos de otra.

\textbf{Clasificación}

\begin{itemize}
  \item \textbf{Hipervisor Tipo 1 (bare-metal):} se instala directamente sobre el hardware físico, sin sistema operativo anfitrión. Ejemplos: VMware ESXi, Microsoft Hyper-V, Xen y KVM. Es el modelo de los centros de datos y de la nube.
  \item \textbf{Hipervisor Tipo 2 (hosted):} se ejecuta como una aplicación sobre un sistema operativo anfitrión. Ejemplos: Oracle VirtualBox y VMware Workstation. Es el modelo usado en este laboratorio para probar varias distribuciones sobre el mismo equipo.
\end{itemize}

\textbf{Características}

\begin{itemize}
  \item Aislamiento: cada VM es independiente; un fallo en una no afecta a las demás.
  \item Partición de recursos: el hipervisor divide CPU, memoria y disco entre las VM.
  \item Encapsulación: la VM completa se representa como archivos (\texttt{.vmx}, \texttt{.vmdk}), portables y copiables.
  \item Independencia del hardware: la VM puede migrar entre anfitriones compatibles.
  \item Instantáneas (snapshots) y, en Tipo 1, migración en vivo sin interrupción del servicio.
\end{itemize}

\textbf{Arquitectura}

Cada VM ejecuta un sistema operativo invitado que cree estar sobre hardware real: CPU virtual (apoyada en las extensiones de virtualización VT-x/AMD-V), memoria virtual, disco virtual y tarjeta de red virtual. El hipervisor intercepta las instrucciones privilegiadas del invitado y las ejecuta en modo supervisor, o las emula mediante dispositivos virtuales; en un Tipo 2, un proceso del anfitrión atiende además las operaciones de entrada/salida (disco y red).

\textbf{Cloud computing}

El cómputo en la nube entrega servicios de procesamiento, almacenamiento, redes, bases de datos y software a través de internet, bajo demanda y con pago por uso. Se distinguen tres modelos de servicio: IaaS (infraestructura como servicio: máquinas virtuales, redes y almacenamiento), PaaS (plataforma para desplegar aplicaciones) y SaaS (software completo). El despliegue puede ser público, privado o híbrido.

\textbf{¿Se usan hipervisores en la nube? Sí.}

La infraestructura de los proveedores de nube pública (AWS, Azure, Google Cloud) está construida sobre hipervisores (KVM, Xen, Hyper-V): el modelo IaaS expone directamente máquinas virtuales; la multitenencia (varios clientes sobre el mismo hardware físico) y la elasticidad (crear o liberar máquinas en segundos) son posibles gracias a la virtualización.

\textbf{Comparación de costos: servidor físico vs.\ servidor en la nube}

\begin{itemize}
  \item Servidor físico: alta inversión inicial (CAPEX) en hardware, más costos fijos de energía, enfriamiento, espacio, mantenimiento y personal; suele quedar sobredimensionado para cubrir los picos de carga.
  \item Servidor en la nube: gasto operativo (OPEX) bajo demanda, sin mantenimiento de hardware y con elasticidad real, pero con costo continuo: en usos 24/7 constantes a largo plazo puede superar el costo del hardware propio.
  \item Conclusión: la nube conviene cuando la carga es variable o el proyecto es nuevo; el hardware propio conviene con carga constante y alta utilización, o cuando se requiere control total sobre el equipo.
\end{itemize}

\textbf{Contenedores}

Un contenedor es una forma de virtualización a nivel del sistema operativo: varios contenedores comparten el kernel del anfitrión, pero sus procesos, sistema de archivos y red están aislados mediante namespaces y cgroups. El contenedor empaqueta la aplicación con todas sus dependencias en una imagen portable y liviana. Ejemplos: Docker, Podman y containerd; la orquestación se hace con Kubernetes.

\textbf{Similitudes y diferencias entre máquinas virtuales y contenedores}

\begin{table}[H]
  \centering
  \begin{tabular}{@{}lll@{}}
    \toprule
    \textbf{Aspecto} & \textbf{Máquinas virtuales} & \textbf{Contenedores} \\
    \midrule
    Qué virtualizan & el hardware & el sistema operativo \\
    Kernel & uno por VM (invitado) & compartido con el anfitrión \\
    Aislamiento & completo y fuerte & a nivel de procesos \\
    Tamaño y arranque & gigabytes, minutos & megabytes, segundos \\
    Densidad por servidor & decenas & cientos o miles \\
    Caso de uso & consolidación, SO heterogéneos & microservicios, CI/CD \\
    \bottomrule
  \end{tabular}
  \caption{Máquinas virtuales vs.\ contenedores.}
  \label{tab:vm-vs-contenedores}
\end{table}

\textbf{Similitudes:} ambos aíslan aplicaciones y recursos, encapsulan la aplicación con su entorno y permiten densidades de despliegue mucho mayores que las físicas. En la práctica se complementan: los contenedores corren sobre las máquinas virtuales en lugar de competir con ellas.

\subsection{Sistemas operativos}

\begin{itemize}
  \item \textbf{Sistema de archivos:} estructura con la que el sistema organiza y almacena los datos en disco: jerarquía de directorios desde la raíz \texttt{/}, archivos con metadatos (permisos, fechas, propietario). En el laboratorio se usaron ext4 (Slackware, con journaling), ZFS (Solaris, con pools, instantáneas y sumas de verificación) y NTFS (Windows).
  \item \textbf{Permisos:} control de acceso a archivos y directorios: tripleta r/w/x para propietario, grupo y otros, con equivalente octal (4/2/1); en Windows se usan ACL (NTFS). Ver sección 3.3.3.
  \item \textbf{Logs:} archivos donde el sistema y los servicios registran eventos: syslog en Unix, Visor de eventos en Windows. Ver sección 3.3.2.
  \item \textbf{Shell:} intérprete de comandos que ejecuta programas y da acceso a los servicios del sistema: interactivo y de scripting en Unix; \texttt{cmd} y PowerShell en Windows. Ver sección 3.7.
\end{itemize}

\section{Desarrollo del tema}

\subsection{Software de virtualización (preguntas del punto 1 del lab)}

\begin{itemize}
  \item ¿Qué son los hipervisores?
  \item ¿Cómo se clasifican?
  \item ¿Cuáles son sus características?
  \item Explicar su arquitectura.
  \item ¿Qué es cloud computing?
  \item ¿Se usan hipervisores en cloud computing? Justificar.
  \item ¿Cuál es la diferencia de costo entre un servidor físico y uno en la nube?
  \item ¿Qué son los contenedores? Explicar su arquitectura.
  \item Similitudes y diferencias entre máquinas virtuales y contenedores.
\end{itemize}

\begin{quote}
\emph{Nota:} este tema se presenta en un video de hasta 5 min (participación de todos los integrantes). El guion/borrador puede ir en esta sección o en un anexo.
\end{quote}

\subsection{Configuración de servidores Unix: Slackware y Solaris}

\subsubsection{Instalación y configuración (VMware, modo experto, sin entorno gráfico)}

\paragraph*{Slackware 15.0}
Instalación en modo consola, sin entorno gráfico, proceso paso a paso. Evidencias del proceso:

\paragraph*{Fase 1 --- Arranque del medio e intento inicial de instalación desde disquete}

\begin{figure}[H]
  \centering
  \includegraphics[width=0.9\linewidth]{00-estado-inicial-instalador.png}
  \caption{Consola de la VM tras arrancar el medio de instalación: la carga del teclado (\texttt{loadkeys es-latam}) y el inicio de \texttt{rc.sshd} fallan porque aún no hay sistema instalado.}
  \label{fig:slack-00}
\end{figure}

\begin{figure}[H]
  \centering
  \includegraphics[width=0.9\linewidth]{01-inicio-instalacion.png}
  \caption{Primer intento de lanzar la instalación desde el disquete: \texttt{sh /mnt/install.sh} responde \texttt{No such file or directory}.}
  \label{fig:slack-01}
\end{figure}

Diagnóstico, montaje y lectura del contenido del disquete de instalación:

\begin{figure}[H]
  \centering
  \includegraphics[width=0.9\linewidth]{01b-diagnostico-floppy.png}
  \caption{Diagnóstico del disquete: listado de \texttt{/dev/fd*} (solo existe \texttt{/dev/fd0}), intento de montaje y comprobación de la red con \texttt{ifconfig eth0 | grep inet}.}
  \label{fig:slack-01b}
\end{figure}

\begin{figure}[H]
  \centering
  \includegraphics[width=0.9\linewidth]{01c-mount-floppy.png}
  \caption{Montaje del disquete: error \texttt{MOUNT\_FAILED} porque \texttt{/dev/fd0} ya estaba montado en \texttt{/mnt}; también se ven las direcciones IPv6 asignadas a la interfaz.}
  \label{fig:slack-01c}
\end{figure}

\begin{figure}[H]
  \centering
  \includegraphics[width=0.9\linewidth]{01d-contenido-floppy.png}
  \caption{Lectura del contenido del disquete con \texttt{ls -la /mnt}: aparece el script de instalación \texttt{INSTALL.SH} (7168 bytes).}
  \label{fig:slack-01d}
\end{figure}

\begin{figure}[H]
  \centering
  \includegraphics[width=0.9\linewidth]{02-instalacion-arrancando.png}
  \caption{Instalación en marcha: montaje del ISO (solo lectura) e instalación de la serie \texttt{a} (129 paquetes, comenzando por \texttt{aaa\_base}).}
  \label{fig:slack-02}
\end{figure}

\paragraph*{Fase 2 --- Instalación de paquetes desde el ISO}

Instalación de paquetes y del parche de instalación:

\begin{figure}[H]
  \centering
  \includegraphics[width=0.9\linewidth]{03-instalando-paquetes.png}
  \caption{Primera pasada fallida: los paquetes del ISO se reportan corruptos (\texttt{tar archive is corrupt, tar returned error code 2}).}
  \label{fig:slack-03}
\end{figure}

\begin{figure}[H]
  \centering
  \includegraphics[width=0.9\linewidth]{03b-patch-install.png}
  \caption{Parche del script de instalación: edición de \texttt{INSTALL.SH} en el disquete (\texttt{sed} para fijar \texttt{installpkg ROOT=/mnt}) y errores de espacio (\texttt{No space left on device}).}
  \label{fig:slack-03b}
\end{figure}

\begin{figure}[H]
  \centering
  \includegraphics[width=0.9\linewidth]{04-reinstalando.png}
  \caption{Segunda pasada de la instalación: \texttt{aaa\_base} y \texttt{aaa\_glibc-solibs} se instalan correctamente (se reinstaló).}
  \label{fig:slack-04}
\end{figure}

\paragraph*{Fase 3 --- Series de paquetes}

Instalación de las series de paquetes:

\begin{figure}[H]
  \centering
  \includegraphics[width=0.9\linewidth]{05-instalando-series.png}
  \caption{Instalación de las series de paquetes (1): paquetes \texttt{python2} y \texttt{python3} de la serie \texttt{l}.}
  \label{fig:slack-05}
\end{figure}

\begin{figure}[H]
  \centering
  \includegraphics[width=0.9\linewidth]{06-instalando-series2.png}
  \caption{Instalación de las series de paquetes (2): librerías \texttt{libxklavier} y \texttt{libxml2} (serie \texttt{l}).}
  \label{fig:slack-06}
\end{figure}

\begin{figure}[H]
  \centering
  \includegraphics[width=0.9\linewidth]{07-instalando-series3.png}
  \caption{Instalación de las series de paquetes (3): firmware \texttt{zd1211} y configuración final del sistema (advertencia de contraseña débil para \texttt{root}).}
  \label{fig:slack-07}
\end{figure}

\paragraph*{Fase 4 --- Primer arranque y login}

Arranque del sistema instalado:

\begin{figure}[H]
  \centering
  \includegraphics[width=0.9\linewidth]{09-boot-con-pae.png}
  \caption{Arranque con el kernel PAE habilitado: generación de claves SSH, inicio de servicios y prompt de login de Slackware 15.0 (\texttt{Welcome to Linux 5.15.19-smp i686}).}
  \label{fig:slack-09-pae}
\end{figure}

\begin{figure}[H]
  \centering
  \includegraphics[width=0.9\linewidth]{09-primer-boot.png}
  \caption{Primer intento de arranque del sistema instalado: LILO 24.2 no carga el kernel porque la CPU no tiene soporte PAE (\texttt{Unable to boot}).}
  \label{fig:slack-09-boot}
\end{figure}

Pantallas de login (HD, resolución 1024 y login como root):

\begin{figure}[H]
  \centering
  \includegraphics[width=0.9\linewidth]{10-login-hd.png}
  \caption{Login en resolución HD (800$\times$600): red inicializada por DHCP (10.0.2.15) y prompt de login.}
  \label{fig:slack-10}
\end{figure}

\begin{figure}[H]
  \centering
  \includegraphics[width=0.9\linewidth]{11-login-1024.png}
  \caption{Login en resolución 1024: arranque completo con \texttt{fsck} de los filesystems y prompt de login.}
  \label{fig:slack-11}
\end{figure}

\begin{figure}[H]
  \centering
  \includegraphics[width=0.9\linewidth]{12-login-root.png}
  \caption{Login como \texttt{root} en la consola (el sistema registra un intento de login fallido previo).}
  \label{fig:slack-12}
\end{figure}

\begin{figure}[H]
  \centering
  \includegraphics[width=0.9\linewidth]{13-sistema.png}
  \caption{Sistema funcionando en consola: \texttt{uname -a}, versión Slackware 15.0, \texttt{df -h} y \texttt{free -m}.}
  \label{fig:slack-13}
\end{figure}

\paragraph*{Fase 5 --- Red y verificación final}

Red y verificación del sistema:

\begin{figure}[H]
  \centering
  \includegraphics[width=0.9\linewidth]{14-red-dhcp.png}
  \caption{Red configurada por DHCP: \texttt{ip addr show eth0} (10.0.2.15/24), ruta por defecto y \texttt{/etc/resolv.conf}.}
  \label{fig:slack-14}
\end{figure}

\begin{figure}[H]
  \centering
  \includegraphics[width=0.9\linewidth]{15-ping-tests.png}
  \caption{Pruebas de conectividad: \texttt{ping} a 8.8.8.8 y a \texttt{www.google.com} con 0\,\% de pérdida de paquetes.}
  \label{fig:slack-15}
\end{figure}

\begin{figure}[H]
  \centering
  \includegraphics[width=0.9\linewidth]{16-directorios.png}
  \caption{Estructura de directorios de la raíz (\texttt{ls -la /}): \texttt{bin}, \texttt{dev}, \texttt{etc}, \texttt{lib}, \texttt{mnt}, \texttt{opt}, \texttt{root}, \texttt{run}, \texttt{srv}, \texttt{sys}, \texttt{tmpmount}, \texttt{usr} y \texttt{var}.}
  \label{fig:slack-16}
\end{figure}

\begin{figure}[H]
  \centering
  \includegraphics[width=0.9\linewidth]{17-teclado-es.png}
  \caption{Cambio de teclado a español con \texttt{loadkeys es} desde la consola.}
  \label{fig:slack-17}
\end{figure}

\paragraph*{Solaris 11.4}
Instalación con el Text Installer, en modo texto, proceso paso a paso. Evidencias del proceso:

\paragraph*{Fase 1 --- Arranque del medio e identidad del sistema}

\begin{figure}[H]
  \centering
  \includegraphics[width=0.9\linewidth]{solaris-02-language.png}
  \caption{Arranque del medio de instalación de Solaris 11.4: selección del layout de teclado (\texttt{Spanish}).}
  \label{fig:sol-02}
\end{figure}

\begin{figure}[H]
  \centering
  \includegraphics[width=0.9\linewidth]{solaris-03-menu-instalacion.png}
  \caption{Menú principal del instalador en modo texto (\texttt{Install Oracle Solaris}, \texttt{Additional Drivers}, \texttt{Terminal type}, \texttt{Reboot}) con una alerta de diagnóstico del kernel (\texttt{SUNW-MSG-ID}).}
  \label{fig:sol-03}
\end{figure}

\begin{figure}[H]
  \centering
  \includegraphics[width=0.9\linewidth]{solaris-04-hostname.png}
  \caption{Definición del nombre del sistema en \texttt{System Identity}: \texttt{solaris}.}
  \label{fig:sol-04}
\end{figure}

Configuración de red y método de configuración:

\begin{figure}[H]
  \centering
  \includegraphics[width=0.9\linewidth]{solaris-05-network.png}
  \caption{Selección de la interfaz de red a configurar (\texttt{net0}, \texttt{e1000g0}) o \texttt{No network}.}
  \label{fig:sol-05}
\end{figure}

\begin{figure}[H]
  \centering
  \includegraphics[width=0.9\linewidth]{solaris-05b-network-method.png}
  \caption{Método de configuración de la interfaz: \texttt{Static} (IP fija) o \texttt{DHCP}.}
  \label{fig:sol-05b}
\end{figure}

\begin{figure}[H]
  \centering
  \includegraphics[width=0.9\linewidth]{solaris-06-disks.png}
  \caption{Selección del disco de destino: disco SATA de 25,0 GB, uso del disco completo o de una partición GPT.}
  \label{fig:sol-06}
\end{figure}

\paragraph*{Fase 2 --- Zona horaria, localización y teclado}

Zona horaria: región América, ciudad Bogotá (\texttt{America/Bogota}):

\begin{figure}[H]
  \centering
  \includegraphics[width=0.9\linewidth]{solaris-07a-timezone-region.png}
  \caption{Zona horaria: selección de la región (\texttt{Americas}).}
  \label{fig:sol-07a}
\end{figure}

\begin{figure}[H]
  \centering
  \includegraphics[width=0.9\linewidth]{solaris-07b-timezone-bogota.png}
  \caption{Zona horaria: selección de la ciudad (\texttt{America/Bogota}).}
  \label{fig:sol-07b}
\end{figure}

Locale en español, territorio, fecha/hora y teclado Latin-American:

\begin{figure}[H]
  \centering
  \includegraphics[width=0.9\linewidth]{solaris-07c-locale.png}
  \caption{Locale: selección del idioma de soporte del sistema (\texttt{Spanish}).}
  \label{fig:sol-07c}
\end{figure}

\begin{figure}[H]
  \centering
  \includegraphics[width=0.9\linewidth]{solaris-07d-territory.png}
  \caption{Territorio del locale: España (\texttt{es\_ES.ISO8859-1}, \texttt{es\_ES.ISO8859-15}, \texttt{es\_ES.UTF-8}).}
  \label{fig:sol-07d}
\end{figure}

\begin{figure}[H]
  \centering
  \includegraphics[width=0.9\linewidth]{solaris-07e-datetime.png}
  \caption{Configuración de fecha y hora del sistema (formato 24\,h, hora UTC).}
  \label{fig:sol-07e}
\end{figure}

\begin{figure}[H]
  \centering
  \includegraphics[width=0.9\linewidth]{solaris-07f-keyboard.png}
  \caption{Selección del teclado: \texttt{Latin-American}.}
  \label{fig:sol-07f}
\end{figure}

\paragraph*{Fase 3 --- Credenciales y resumen de configuración}

Credenciales del usuario \texttt{root} y opción de registro de soporte:

\begin{figure}[H]
  \centering
  \includegraphics[width=0.9\linewidth]{solaris-08-users.png}
  \caption{Credenciales del sistema: contraseña de \texttt{root} y creación opcional de una cuenta de usuario.}
  \label{fig:sol-08}
\end{figure}

\begin{figure}[H]
  \centering
  \includegraphics[width=0.9\linewidth]{solaris-08b-support-registration.png}
  \caption{Opción de registro de soporte de My Oracle Support (telemetría opcional).}
  \label{fig:sol-08b}
\end{figure}

\begin{figure}[H]
  \centering
  \includegraphics[width=0.9\linewidth]{solaris-09-summary.png}
  \caption{Resumen de configuración antes de instalar: Solaris 11.4 x86, disco 25,0 GB, red por DHCP, \texttt{America/Bogota}, español y teclado Latin-American.}
  \label{fig:sol-09}
\end{figure}

\begin{figure}[H]
  \centering
  \includegraphics[width=0.9\linewidth]{solaris-10-instalando.png}
  \caption{Instalación de Oracle Solaris en curso (inicialización del destino completada).}
  \label{fig:sol-10}
\end{figure}

\begin{figure}[H]
  \centering
  \includegraphics[width=0.9\linewidth]{solaris-11-completado.png}
  \caption{Instalación finalizada: el sistema indica reiniciar para arrancar Solaris 11.4.}
  \label{fig:sol-11}
\end{figure}

\begin{figure}[H]
  \centering
  \includegraphics[width=0.9\linewidth]{solaris-12-login.png}
  \caption{Primer arranque: SunOS 5.11 (Solaris 11.4), hostname \texttt{solaris} y prompt de login de consola.}
  \label{fig:sol-12}
\end{figure}

\begin{figure}[H]
  \centering
  \includegraphics[width=0.9\linewidth]{solaris-13-sistema.png}
  \caption{Login como \texttt{root} en la consola de Solaris 11.4 (\texttt{Oracle Corporation SunOS 5.11}).}
  \label{fig:sol-13}
\end{figure}

\paragraph*{Fase 4 --- Instalación, primer arranque y login}

Otros temas de la sección:

\begin{itemize}
  \item Paquetes instalados: solo los necesarios para operación básica y conectividad de red.
  \item Archivos generados durante la instalación por el software de virtualización y su propósito (\texttt{.vmx}, \texttt{.vmdk}, \texttt{.nvram}, \texttt{.vmsd}, \texttt{.vmxf}, logs\ldots).
  \item ¿Es posible convertir una VM de VMware a VirtualBox y viceversa? (OVF/OVA, \texttt{vmware-vdiskmanager}, \texttt{VBoxManage clonehd}\ldots)
\end{itemize}

\subsubsection{Usuarios y grupos}

En cada sistema Unix (Slackware y Solaris) se crearon 4 usuarios con una descripción significativa (campo \texttt{-c}), con su directorio home dentro de \texttt{/usuarios} en la raíz del filesystem y perteneciendo a \emph{un solo grupo}: los usuarios 1 y 2 (\texttt{claudia}, \texttt{john}) al grupo Accounting, y los usuarios 3 y 4 (\texttt{fabian}, \texttt{diego}) al grupo IT.

\paragraph*{Slackware 15.0}

En Slackware los nombres de grupo no pueden contener mayúsculas: el \texttt{shadow-utils} de la distribución rechaza \texttt{groupadd Accounting} con \texttt{'Accounting' is not a valid group name}. Por eso los grupos se crearon en minúsculas: \texttt{accounting} e \texttt{it}.

Comandos ejecutados (como \texttt{root}, o con \texttt{sudo} desde el usuario \texttt{vagrant}):

\begin{lstlisting}
# Directorio base de los homes
mkdir -p /usuarios

# Grupos (en minusculas por la limitacion de shadow-utils)
groupadd accounting
groupadd it

# Usuarios 1 y 2 -> solo accounting
useradd -d /usuarios/claudia -m -c "Claudia - usuario con el nombre de la profesora" -g accounting claudia
useradd -d /usuarios/john    -m -c "John - usuario con el nombre del profesor"     -g accounting john

# Usuarios 3 y 4 -> solo it
useradd -d /usuarios/fabian -m -c "Fabian - usuario con el nombre del profesor" -g it fabian
useradd -d /usuarios/diego  -m -c "Diego - usuario con el nombre del profesor"  -g it diego

# Asignar contrasena a cada usuario
passwd claudia
passwd john
passwd fabian
passwd diego
\end{lstlisting}

Significado de las opciones: \texttt{-d} define la ruta del home; \texttt{-m} lo crea y copia las plantillas de \texttt{/etc/skel}; \texttt{-c "..."} guarda la descripción del usuario (visible en el campo GECOS de \texttt{/etc/passwd}); \texttt{-g} fija el grupo \emph{primario}, con lo que cada usuario queda en ese único grupo (cumple la consigna ``solo Accounting / solo IT'', sin grupos secundarios). Si hubiera que corregir algo posteriormente, se usa \texttt{usermod -g <grupo> <usuario>} para cambiar el grupo primario o \texttt{groupdel <grupo>} para eliminar un grupo vacío.

Resultado verificado (04/08/2026):

\begin{table}[H]
  \centering
  \begin{tabular}{@{}lllll@{}}
    \toprule
    \textbf{Usuario} & \textbf{uid} & \textbf{Grupo primario} & \textbf{gid} & \textbf{Home} \\
    \midrule
    claudia & 1001 & accounting & 1002 & /usuarios/claudia \\
    john    & 1002 & accounting & 1002 & /usuarios/john \\
    fabian  & 1003 & it         & 1003 & /usuarios/fabian \\
    diego   & 1004 & it         & 1003 & /usuarios/diego \\
    \bottomrule
  \end{tabular}
  \caption{Usuarios creados en Slackware 15.0.}
  \label{tab:slack-users}
\end{table}

Verificación con \texttt{id}, \texttt{getent} y \texttt{ls}:

\begin{lstlisting}[language={}]
# id claudia
uid=1001(claudia) gid=1002(accounting) groups=1002(accounting)

# getent group accounting
accounting:x:1002:

# ls -ld /usuarios
drwxr-xr-x 3 root root 4096 /usuarios

# ls -la /usuarios
drwxr-xr-x  2 claudia accounting 4096 ... claudia
\end{lstlisting}

Notas de la ejecución:

\begin{itemize}
  \item El acceso a la VM es con el usuario \texttt{vagrant} (SSH o consola). Desde esa cuenta, el PATH \emph{no} incluye \texttt{/usr/sbin} (Slackware solo lo agrega al PATH de root), por lo que \texttt{sudo groupadd} falla con \texttt{sudo: groupadd: command not found}. Se resolvió logueándose como \texttt{root} o invocando la ruta completa: \texttt{sudo /usr/sbin/groupadd it}, \texttt{sudo /usr/sbin/useradd \ldots}
  \item \texttt{groupadd accounting} puede responder ``already exists'' (el grupo existía de una prueba anterior); no es un error, el grupo queda correcto.
\end{itemize}

\paragraph*{Solaris 11.4}

En Solaris los nombres de grupo sí admiten mayúsculas, por lo que los grupos se crearon como pide el lab: \texttt{Accounting} (GID 100) e \texttt{IT} (GID 101). Los nombres de usuario van en minúsculas.

Comandos ejecutados (como \texttt{root} en la consola de la VM):

\begin{lstlisting}
# Directorio base de los homes
mkdir -p /usuarios

# Grupos
groupadd Accounting
groupadd IT

# Usuarios 1 y 2 -> solo Accounting
useradd -d /usuarios/claudia -m -c "Claudia - usuario con el nombre de la profesora" -g Accounting -s /usr/bin/bash claudia
useradd -d /usuarios/john    -m -c "John - usuario con el nombre del profesor"     -g Accounting -s /usr/bin/bash john

# Usuarios 3 y 4 -> solo IT
useradd -d /usuarios/fabian -m -c "Fabian - usuario con el nombre del profesor" -g IT -s /usr/bin/bash fabian
useradd -d /usuarios/diego  -m -c "Diego - usuario con el nombre del profesor"  -g IT -s /usr/bin/bash diego

# Asignar contrasena a cada usuario
passwd claudia
passwd john
passwd fabian
passwd diego
\end{lstlisting}

Mismas opciones que en Slackware, con dos diferencias de sintaxis: \texttt{-s /usr/bin/bash} define el shell de login (por defecto Solaris usa \texttt{/usr/bin/sh}), y el directorio padre \texttt{/usuarios} debe existir antes de \texttt{useradd -m} (si falta, \texttt{useradd} falla con \texttt{UX: useradd: ERROR: \ldots /usuarios: No such file or directory}; el \texttt{mkdir -p} del paso 1 lo resuelve).

Resultado verificado (04/08/2026):

\begin{table}[H]
  \centering
  \begin{tabular}{@{}lllll@{}}
    \toprule
    \textbf{Usuario} & \textbf{uid} & \textbf{Grupo primario} & \textbf{gid} & \textbf{Home} \\
    \midrule
    claudia & 100 & Accounting & 100 & /usuarios/claudia \\
    john    & 101 & Accounting & 100 & /usuarios/john \\
    fabian  & 102 & IT         & 101 & /usuarios/fabian \\
    diego   & 103 & IT         & 101 & /usuarios/diego \\
    \bottomrule
  \end{tabular}
  \caption{Usuarios creados en Solaris 11.4.}
  \label{tab:sol-users}
\end{table}

La verificación se hizo con los mismos comandos (\texttt{id <usuario>}, \texttt{getent group Accounting IT}, \texttt{ls -la /usuarios}), comprobando que cada home pertenece a su usuario y a su grupo, y que cada usuario aparece en un solo grupo.

\paragraph*{Diferencia entre ambos sistemas}

Más allá de la limitación de mayúsculas de Slackware, los UID/GID de Solaris comienzan en 100 (numeración del instalador), mientras que Slackware los asigna desde 1000/1001; en el resto, la sintaxis y el resultado final son equivalentes.

\subsubsection{Sistema de archivos}

\begin{itemize}
  \item ¿Qué es el sistema de archivos?
  \item ¿Cuál se usó en la instalación? \textbf{Slackware 15.0: ext4}; \textbf{Solaris 11.4: ZFS} (el instalador de Solaris crea el pool \texttt{rpool} sobre el disco de destino por defecto).
  \item Características de cada uno.
\end{itemize}

\subsubsection{Red}

\begin{itemize}
  \item En el hogar la red se configuró automáticamente por DHCP en modo bridge: en Slackware quedó registrada la configuración DHCP (Figura~\ref{fig:slack-14}) y se ejecutaron pruebas de ping (Figura~\ref{fig:slack-15}). La red de la universidad, en cambio, \textbf{no tiene DHCP}: es una red estática (10.2.0.0/16, gateway y DNS 10.2.65.1, dominio \texttt{is.escuelaing.edu.co}) en la que cada equipo debe configurarse a mano.
  \item ¿Qué significan ``Bridge Mode'' y ``NAT Mode''?
    \begin{itemize}
      \item \textbf{Bridge:} la VM comparte la tarjeta de red física y aparece en la LAN como un equipo más, con su propia dirección IP; los demás equipos la ven directamente.
      \item \textbf{NAT:} la VM sale a la red a través del anfitrión, que traduce las direcciones; la VM queda oculta detrás del host y solo tiene salida, no es alcanzable desde fuera.
    \end{itemize}
  \item Configuración manual (estática):
    \begin{itemize}
      \item IP: \texttt{10.2.78.n} (rango asignado por el instructor; sin duplicar entre sistemas) $\rightarrow$ Slackware=\texttt{10.2.78.64}, Windows Core=\texttt{10.2.78.66}, Windows GUI=\texttt{10.2.78.67}, Android=\texttt{10.2.78.68}, host=\texttt{10.2.78.69}; Solaris=\texttt{10.2.78.65}
      \item Máscara: \texttt{255.255.0.0}
      \item Gateway: \texttt{10.2.65.1}
      \item DNS: \texttt{10.2.65.1}
    \end{itemize}
\end{itemize}

\textbf{Configuración estática en Slackware 15.0 (editada con nano)}

En Slackware la configuración de red se edita en \texttt{/etc/rc.d/rc.inet1.conf} con el editor \texttt{nano}.

\begin{enumerate}
  \item Abrir el archivo de red:
  \begin{verbatim}
nano /etc/rc.d/rc.inet1.conf
  \end{verbatim}
  \item Dentro de nano, con las flechas, se modificaron 4 valores (guardar con \texttt{Ctrl+O}, salir con \texttt{Ctrl+X}):
  \begin{center}
  \begin{tabular}{lll}
    \toprule
    \textbf{Línea} & \textbf{Valor} \\
    \midrule
    \texttt{USE\_DHCP[0]="yes"} & \texttt{USE\_DHCP[0]="no"} \\
    \texttt{IPADDR[0]=""} & \texttt{IPADDR[0]="10.2.78.64"} \\
    \texttt{NETMASK[0]=""} & \texttt{NETMASK[0]="255.255.0.0"} \\
    \texttt{GATEWAY=""} (agregada al final) & \texttt{GATEWAY="10.2.65.1"} \\
    \bottomrule
  \end{tabular}
  \end{center}
  \item El DNS se configuró en \texttt{/etc/resolv.conf} (también con nano). En Slackware este archivo no es un enlace a resolvconf, por lo que la edición directa persiste: \texttt{nameserver 10.2.65.1}.
  \item Aplicar y verificar:
  \begin{verbatim}
/etc/rc.d/rc.inet1 restart
ip addr show
ip route show
  \end{verbatim}
  Resultado verificado: \texttt{eth0} con \texttt{10.2.78.64/16} y ruta por defecto \texttt{default via 10.2.65.1}.
  \item Para alternar entre la red de la universidad (estática) y la del hogar (DHCP/NAT), se dejaron dos scripts dentro del sistema: \texttt{/usr/local/bin/uni} (estática) y \texttt{/usr/local/bin/casa} (DHCP).
\end{enumerate}

\textbf{Configuración estática en Solaris 11.4 (comandos \texttt{ipadm})}

Solaris 11 usa el modelo SMF/\texttt{ipadm} (no archivos de script como \texttt{rc.inet1.conf}); la IP estática se verificó con \texttt{ipadm show-addr}.

\begin{enumerate}
  \item Ver y eliminar la configuración DHCP existente:
  \begin{verbatim}
ipadm show-addr
ipadm delete-addr net0/v4
  \end{verbatim}
  \item Crear la dirección estática y el alias de red NAT:
  \begin{verbatim}
ipadm create-addr -T static -a local=10.2.78.65/16 net0/v4
ipadm create-addr -T static -a local=10.0.2.15/24 net0/v4nat
  \end{verbatim}
  \item Ruta por defecto persistente y DNS:
  \begin{verbatim}
route -p add default 10.2.65.1
echo "nameserver 10.2.65.1" > /etc/resolv.conf
  \end{verbatim}
  \item Verificación: \texttt{ipadm show-addr} muestra \texttt{net0/v4} = \texttt{10.2.78.65/16} estática.
  \item Scripts de alternancia uni/casa: \texttt{/root/uni.sh} y \texttt{/root/casa.sh}.
\end{enumerate}

\paragraph*{Pruebas de conectividad en la red de la universidad}

Las pruebas se ejecutaron en la red del laboratorio con las direcciones estáticas de la sección anterior; los resultados se resumen a continuación:

\begin{itemize}
  \item \textbf{Host (portátil físico):} \texttt{10.2.78.69/16} con gateway y DNS \texttt{10.2.65.1}, configurado con \texttt{netsh} sobre el adaptador Ethernet.
  \item \textbf{Slackware (10.2.78.64):} estática en \texttt{/etc/rc.d/rc.inet1.conf} y DNS en \texttt{/etc/resolv.conf}; la verificación de la interfaz y de las rutas confirmó eth0 con \texttt{10.2.78.64/16} y ruta por defecto \texttt{default via 10.2.65.1} (la interfaz conserva además el alias NAT \texttt{10.0.2.15/24} agregado en \texttt{rc.local} para el escenario del hogar).
  \item \textbf{Windows Server Core (10.2.78.66):} estática aplicada con \texttt{netsh}. Como Windows bloquea el ping entrante por defecto, se habilitó el ICMP en el firewall (regla \texttt{netsh advfirewall} para \texttt{icmpv4:8}); la IP se verificó con \texttt{ipconfig} (hostname \texttt{WIN-O95R63RBHK9}).
  \item \textbf{Windows Server GUI (10.2.78.67):} estática aplicada con \texttt{netsh} (hostname \texttt{WIN-3VMPCF0KCDN}).
  \item \textbf{Android (10.2.78.68):} IP estática sobre la interfaz ethernet \texttt{wifi\_eth} de android-x86. Android no ofrece la configuración estática estándar desde la consola de emergencia, por lo que se usó el \emph{policy routing} de Linux con una tabla propia (tabla 97):
\begin{verbatim}
ip rule add from all lookup main pref 50
ip route add default via 10.2.65.1 dev wifi_eth table 97
ip route add 10.2.0.0/16 dev wifi_eth src 10.2.78.68 table 97
\end{verbatim}
  \textbf{Limitación documentada:} la resolución DNS por nombre no funciona desde la consola de emergencia (Android sin el servicio \texttt{netd} no consulta el DNS configurado): el ping por IP funciona, el ping por nombre no resuelve.
  \item \textbf{Solaris (10.2.78.65):} instalado y probado; IP estática con \texttt{ipadm}, pings OK.
\end{itemize}

Tabla de resultados reales (0\,\% de pérdida de paquetes en todas las pruebas OK):

\begin{table}[H]
  \centering
  \footnotesize
  \begin{tabularx}{\linewidth}{@{}p{2.9cm}XXXXX@{}}
    \toprule
    \textbf{Prueba} & \textbf{Slackware .64} & \textbf{Core .66} & \textbf{GUI .67} & \textbf{Android .68} & \textbf{Solaris .65} \\
    \midrule
    IP propia & \texttt{10.2.78.64/16} & \texttt{10.2.78.66/16} & \texttt{10.2.78.67/16} & \texttt{10.2.78.68/16} (\texttt{wifi\_eth}) & \texttt{10.2.78.65/16} \\
    Ping a la propia IP & OK --- 0\,\% pérdida & --- & --- & --- & OK --- 0\,\% pérdida \\
    Ping al gateway (10.2.65.1) & OK --- 0\,\% pérdida & OK --- 0\,\% pérdida & OK --- 0\,\% pérdida & OK & OK --- 0\,\% pérdida \\
    Ping externo (8.8.8.8) & OK --- 0\,\% pérdida & OK --- 0\,\% pérdida & OK --- 0\,\% pérdida & OK & OK --- 0\,\% pérdida \\
    Ping inter-VM & OK a Core (.66) & OK a Slackware (.64) y GUI (.67) & OK a Slackware (.64) & OK a Slackware (.64) & OK --- 0\,\% pérdida \\
    Resolución DNS (\texttt{ping www.google.com}) & OK --- 0\,\% pérdida & OK --- 0\,\% pérdida & OK --- 0\,\% pérdida & No resuelve por nombre (limitación) & OK \\
    \bottomrule
  \end{tabularx}
  \caption{Resultados reales de las pruebas de conectividad en la red de la universidad.}
  \label{tab:ping}
\end{table}

Latencia media medida desde Slackware: 1,1\,ms al gateway, 5,2\,ms a 8.8.8.8 y 4,4\,ms a \texttt{www.google.com} (RTT máximo 6\,ms).

\begin{figure}[H]
  \centering
  \includegraphics[width=0.85\linewidth]{red/android-pings-uni2.png}
  \caption{Pruebas de ping desde Android-x86 (continuación): conectividad externa e inter-VM.}
  \label{fig:and-ping-2}
\end{figure}

Pruebas OK (0\,\% de pérdida de paquetes): gateway, 8.8.8.8, \texttt{www.google.com} e inter-VM entre todos los sistemas.

\subsubsection{Comparación de experiencias de instalación}

Ambos sistemas se instalaron en modo texto, sin entorno gráfico, sobre VMware. La experiencia fue muy distinta:

\begin{table}[H]
  \centering
  \footnotesize
  \begin{tabularx}{\linewidth}{@{}p{2.6cm}XX@{}}
    \toprule
    \textbf{Aspecto} & \textbf{Slackware 15.0} & \textbf{Solaris 11.4} \\
    \midrule
    Instalador & Script de consola con series de paquetes & Text Installer guiado (menús paso a paso) \\
    Dificultad & Alta: cada decisión (particionado, series, arranque) es manual & Baja-media: el asistente guía todo el proceso \\
    Paquetes & Series (\texttt{a}, \texttt{ap}, \texttt{d}, \texttt{l}, \texttt{n}, \texttt{x}, \texttt{xap}\ldots); sin resolución automática de dependencias (\texttt{installpkg}) & Selección por grupos; IPS (\texttt{pkg}) resuelve dependencias \\
    Tiempo & Largo: varias pasadas y correcciones (parche de \texttt{INSTALL.SH}, reinstalación, problema de arranque por PAE) & Corto: una sola pasada guiada \\
    Control & Total sobre cada etapa & El instalador automatiza (pool ZFS \texttt{rpool}, boot environments) \\
    \bottomrule
  \end{tabularx}
  \caption{Comparación de la experiencia de instalación: Slackware vs.\ Solaris.}
  \label{tab:cmp-instalacion}
\end{table}

En Slackware la instalación fue un ejercicio de administración pura: hubo que diagnosticar el medio de instalación, parchear el script \texttt{INSTALL.SH} (los paquetes del ISO se reportaban corruptos), reinstalar el sistema y resolver un fallo de arranque por falta de soporte PAE (LILO no cargaba el kernel), para finalmente elegir las series de paquetes a instalar. El resultado es un sistema mínimo y totalmente conocido, sin gestor de dependencias: los paquetes se administran con \texttt{installpkg}/\texttt{upgradepkg}/\texttt{removepkg}.

En Solaris la instalación fue mucho más rápida y guiada: el Text Installer conduce paso a paso (idioma, hostname, red, discos, zona horaria, teclado, credenciales y resumen final), crea automáticamente el pool ZFS \texttt{rpool} y deja un sistema con gestión de paquetes IPS y boot environments (instantáneas del sistema de archivos raíz que permiten volver atrás).

Conclusión: Solaris privilegia la velocidad y la automatización; Slackware exige entender cada etapa, lo que lo hace más educativo, aunque más lento de instalar.

\subsection{Comprensión y administración de los sistemas operativos}

\subsubsection{Estructura de directorios}

\begin{itemize}
  \item Listar directorios, describir su contenido y comparar Slackware vs.\ Solaris (evidencias de la estructura raíz de Slackware: Figuras~\ref{fig:est-raiz},~\ref{fig:est-raiz-1} y~\ref{fig:slack-16}). En la raíz de Slackware se ven: \texttt{bin}, \texttt{boot}, \texttt{cdrom}, \texttt{dev}, \texttt{etc}, \texttt{home}, \texttt{lib}, \texttt{lost+found}, \texttt{media}, \texttt{mnt}, \texttt{opt}, \texttt{root}, \texttt{run}, \texttt{sbin}, \texttt{tmp}, \texttt{tmpmount}, \texttt{usr}, \texttt{var} y \texttt{usuarios}.
    \begin{itemize}
      \item \textbf{Archivos de configuración (\texttt{/etc}):} texto plano editable por el administrador: \texttt{/etc/passwd}, \texttt{/etc/shadow}, \texttt{/etc/resolv.conf} y \texttt{/etc/rc.d/rc.inet1.conf} en Slackware; en Solaris la red vive en \texttt{/etc/inet} y los servicios en \texttt{/etc/svc}. Es la ``memoria'' de la configuración del sistema.
      \item \textbf{Ejecutables (\texttt{/bin}, \texttt{/sbin}, \texttt{/usr/bin}, \texttt{/usr/sbin}):} los programas del sistema están separados por función y por orden de montaje: \texttt{/bin} y \texttt{/sbin} contienen lo esencial para arrancar y reparar el sistema antes de montar \texttt{/usr} (en Slackware 15 son enlaces a \texttt{usr/bin} y \texttt{usr/sbin}, aunque la convención se conserva); \texttt{/usr/bin} tiene los programas de uso general y \texttt{/usr/sbin} las herramientas de administración (reservadas a root). Varias ubicaciones = jerarquía histórica: lo mínimo para el arranque vs.\ el resto del sistema.
      \item \textbf{Logs (\texttt{/var/log}):} aquí se registran los eventos del sistema (ver sección 3.3.2): en Slackware \texttt{/var/log/syslog}, \texttt{/var/log/messages} y \texttt{/var/log/secure}; en Solaris \texttt{/var/adm/messages} y \texttt{/var/adm/utmpx}.
      \item \textbf{Montaje de dispositivos externos (\texttt{/mnt}, \texttt{/media}):} \texttt{/mnt} para montajes manuales (el USB de la prueba: \texttt{mount /dev/sdb1 /mnt/usb}) y \texttt{/media} para medios removibles montados automáticamente.
    \end{itemize}
  \item Conectar una USB y hacerla visible en la VM: se emuló una USB en Slackware (disco virtual de 512\,MiB adjuntado a la VM, visible como \texttt{VBOX HARDDISK}); flujo completo en \texttt{lab-evidencias/usb/slackware-usb.txt}:
    \begin{enumerate}
      \item El dispositivo aparece como \texttt{/dev/sdb} (512\,MiB, 1048576 sectores) y se particiona con \texttt{fdisk /dev/sdb}.
      \item Formateo ext4: \texttt{mkfs.ext4 /dev/sdb1} (crea el journal: \texttt{Creating journal (4096 blocks): done}).
      \item Punto de montaje y montaje: \texttt{mkdir -p /mnt/usb}; \texttt{mount /dev/sdb1 /mnt/usb} $\rightarrow$ el dispositivo queda montado en \texttt{/mnt/usb} (488M de capacidad).
      \item Copia de un archivo de prueba: \texttt{cp evidencia.txt /mnt/usb/} $\rightarrow$ aparece \texttt{evidencia.txt} (23 bytes) junto al directorio \texttt{lost+found}.
      \item Desmontaje: \texttt{umount /mnt/usb} $\rightarrow$ ``sdb desmontado OK''.
    \end{enumerate}
    Los puntos de montaje estándar para dispositivos externos en Slackware son \texttt{/mnt} (montajes manuales) y \texttt{/media} (medios removibles).
  \item Diferencias observadas entre ambos sistemas: Slackware presenta la estructura clásica de Linux (\texttt{/boot}, \texttt{/lib}, \texttt{/lost+found}, \texttt{/media}, \texttt{/mnt}, \texttt{/opt}, \texttt{/srv}, \texttt{/tmp}, \texttt{/home}); Solaris añade directorios propios de SunOS (\texttt{/kernel} y \texttt{/platform} con los módulos del kernel, \texttt{/devices} con los nodos de dispositivo gestionados por devfs, \texttt{/etc/inet} para la red y \texttt{/var/adm} para los logs del sistema). Ambos siguen el estándar común de \texttt{/etc}, \texttt{/bin}, \texttt{/usr}, \texttt{/var}, \texttt{/tmp} y \texttt{/root}.
\end{itemize}

\begin{figure}[H]
  \centering
  \includegraphics[width=0.9\linewidth]{slackware-estructura-raiz.png}
  \caption{Estructura raíz de Slackware (\texttt{ls -la /}): \texttt{bin}, \texttt{boot}, \texttt{cdrom}, \texttt{dev}, \texttt{etc}, \texttt{home}, \texttt{lib}, \texttt{lost+found}, \texttt{media}, \texttt{mnt}, \texttt{opt}, \texttt{root}, \texttt{run}, \texttt{sbin}, \texttt{tmp}, \texttt{tmpmount}, \texttt{usuarios} y \texttt{var}.}
  \label{fig:est-raiz}
\end{figure}

\begin{figure}[H]
  \centering
  \includegraphics[width=0.9\linewidth]{slackware-estructura-raiz-1.png}
  \caption{Estructura raíz de Slackware (2): detalle de la consola durante el listado (comando \texttt{ls -} con error de sintaxis).}
  \label{fig:est-raiz-1}
\end{figure}

\subsubsection{Logs y syslog}

\begin{itemize}
  \item ¿Qué son los archivos de log? Archivos donde el sistema y los servicios registran eventos (arranque del kernel, inicios de sesión, errores de red, avisos de aplicaciones), cada entrada con fecha/hora, origen y mensaje.
  \item ¿Qué es syslog? Protocolo y servicio de registro central de eventos. Archivos principales: en Slackware \texttt{/var/log/syslog}, \texttt{/var/log/messages} y \texttt{/var/log/secure} (inicios de sesión y \texttt{sudo}); en Solaris \texttt{/var/adm/messages} (syslog del sistema) y \texttt{/var/adm/utmpx}/\texttt{wtmpx} (sesiones). En Slackware el demonio es \texttt{syslogd}; en Solaris el servicio SMF \texttt{svc:/system/system-log}.
  \item Tipos de información registrada y estructura de una entrada: fecha/hora, hostname, programa (o \texttt{kernel}) con su PID y el mensaje; en Solaris las entradas incluyen además un ID de evento (p.\,ej.\ \texttt{[ID 936769 kern.info]}).
  \item Cinco ejemplos reales de eventos registrados:
    \begin{enumerate}
      \item \textbf{Kernel (Slackware)}: mensajes del kernel con VirtualBox Guest Additions --- \texttt{Aug 5 20:24:19 slackware kernel: \ldots main 7.2.12 r174389 started. Verbose level = 0} (\texttt{lab-evidencias/logs/slackware-syslog.txt}).
      \item \textbf{sshd/sudo (Slackware)}: \texttt{Aug 6 03:00:14 slackware sudo: vagrant : PWD=/home/vagrant ; USER=root ; COMMAND=/usr/bin/tail -6 /var/log/secure}, seguido de \texttt{pam\_unix(sudo:session): session opened for user root} (\texttt{lab-evidencias/logs/slackware-ssh.txt}).
      \item \textbf{Kernel (Solaris)}: detección de dispositivos en el arranque --- \texttt{Aug 5 22:05:09 solaris genunix: [ID 936769 kern.info] pm0 is /pseudo/pm@0} (\texttt{lab-evidencias/logs/solaris-messages.txt}).
      \item \textbf{Red (Solaris)}: enlace de red activo --- \texttt{Aug 5 22:05:22 solaris mac: [ID 435574 kern.info] NOTICE: e1000g0 link up, 1000 Mbps, full duplex} (\texttt{lab-evidencias/logs/solaris-messages.txt}).
      \item \textbf{sendmail (Solaris)}: \texttt{Aug 5 22:09:36 solaris sendmail[902]: [ID 702911 mail.crit] My unqualified host name (solaris) unknown; sleeping for retry} (\texttt{lab-evidencias/logs/solaris-messages.txt}).
    \end{enumerate}
  \item ¿Funciona syslog en los sistemas instalados? \textbf{Sí, en ambos}: en Slackware los eventos quedan en \texttt{/var/log/syslog}, \texttt{/var/log/messages} y \texttt{/var/log/secure} (evidencias \texttt{slackware-syslog.txt} y \texttt{slackware-ssh.txt}); en Solaris en \texttt{/var/adm/messages} (evidencia \texttt{solaris-messages.txt}).
\end{itemize}

\subsubsection{Permisos}

\begin{itemize}
  \item Cómo funcionan los permisos en cada sistema: tripleta r/w/x aplicada a propietario/grupo/otros, visible con \texttt{ls -l} (en Solaris los grupos aparecen como \texttt{Accounting}/\texttt{IT}); los valores rwx se traducen al modo octal (4/2/1). Solaris suma además el modelo ZFS/ACL.
  \item Modificación con representación simbólica (\texttt{chmod g+w}, \texttt{chmod u+x,g-w}) y numérica (\texttt{chmod 750}, \texttt{chmod 640}, \texttt{chmod 770}).
  \item Pruebas de impacto reales en \textbf{Slackware} (evidencia \texttt{lab-evidencias/logs/\allowbreak slackware-permisos.txt}):
    \begin{itemize}
      \item Directorio \texttt{/lab-permisos} propiedad de \texttt{root:it}: \texttt{chmod} numérico (750 sobre el directorio en el proceso y 640 sobre \texttt{archivo.txt}) y simbólico \texttt{g+w}; estado final verificado: \texttt{drwxrwx--- 2 root it} (770).
      \item \textbf{fabian} (grupo \texttt{it}) puede crear archivos: \texttt{touch /lab-permisos/fabian.txt} $\rightarrow$ OK (\texttt{FABIAN\_OK}).
      \item \textbf{claudia} (grupo \texttt{accounting}) no puede: \texttt{touch: cannot touch '/lab-permisos/claudia.txt': Permission denied}.
      \item Archivo \texttt{archivo.txt} con 640 (\texttt{-rw-r----- root:root}): solo root lo lee. Tras \texttt{chgrp it} y permisos simbólicos de lectura/escritura al grupo, \textbf{fabian} logra leerlo y \textbf{claudia} sigue denegada (\texttt{Permission denied}).
      \item Simbólico \texttt{chmod u+x,g-w} sobre el archivo $\rightarrow$ \texttt{-rwxr-----} (ejecución al propietario, escritura quitada al grupo).
    \end{itemize}
  \item Pruebas de impacto reales en \textbf{Solaris} (evidencia \texttt{lab-evidencias/logs/\allowbreak solaris-permisos.txt}):
    \begin{itemize}
      \item Directorio \texttt{/lab-permisos} con \texttt{chmod 770 root:IT} $\rightarrow$ \texttt{drwxrwx--- 2 root IT}.
      \item \textbf{fabian} (grupo \texttt{IT}) crea (\texttt{FABIAN\_OK}) y lee \texttt{archivo.txt} sin errores.
      \item \textbf{claudia} (grupo \texttt{Accounting}) denegada: \texttt{touch: cannot stat /lab-permisos/claudia.txt: Permission denied} y \texttt{cat: cannot open /lab-permisos/archivo.txt: Permission denied}.
      \item Simbólico \texttt{chmod g-w} aplicado para quitar escritura al grupo.
    \end{itemize}
  \item Conclusión de la matriz de impacto: un usuario del grupo propietario del directorio (fabian/IT) puede crear y leer; uno ajeno (claudia/Accounting) recibe \texttt{Permission denied} en ambos sistemas, confirmando que los permisos funcionan igual en Slackware y Solaris.
\end{itemize}

\subsection{Windows Server --- Fase 1 (Server Core, sin interfaz gráfica)}

Instalación de la edición \textbf{Server Core} (sin interfaz gráfica) en VMware. Evidencias:

\subsubsection{Instalación}

\begin{figure}[H]
  \centering
  \includegraphics[width=0.9\linewidth]{windows-01-instalando.png}
  \caption{Instalación de Windows Server 2025 en curso: copia de archivos y preparación del sistema (edición Core).}
  \label{fig:win-01}
\end{figure}

\subsubsection{Primer arranque (SConfig)}

\begin{figure}[H]
  \centering
  \includegraphics[width=0.9\linewidth]{windows-02-sconfig-primer-arranque.png}
  \caption{Primer arranque con SConfig: Windows Server 2025 Standard Evaluation, hostname \texttt{WIN-095R63R3HG}, dominio \texttt{WORKGROUP} y RDP habilitado.}
  \label{fig:win-02}
\end{figure}

Datos verificados de la instalación:

\begin{itemize}
  \item Hostname temporal asignado por el instalador: \texttt{WIN-095R63R3HG}.
  \item Cuenta local \texttt{Administrator} con contraseña \texttt{Server2025!}.
  \item Dominio: \texttt{WORKGROUP} (la VM no se unió a ningún dominio).
\end{itemize}

Red y pruebas (detalle en el párrafo ``Pruebas de conectividad en la red de la universidad'' de la sección 3.2.4):

\begin{itemize}
  \item Red estática aplicada con \texttt{netsh}: 10.2.78.66/16, GW 10.2.65.1, DNS 10.2.65.1.
  \item ICMP habilitado en el firewall (Windows bloquea el ping entrante por defecto): regla \texttt{netsh advfirewall} para \texttt{icmpv4:8}.
  \item Pruebas OK (0\,\% de pérdida): gateway, 8.8.8.8, \texttt{www.google.com} e inter-VM con Slackware (.64) y GUI (.67).
\end{itemize}

Consigna pendiente:

\begin{itemize}
  \item Acceso remoto (RDP): [PENDIENTE: verificar el estado de RDP en la VM].
\end{itemize}

\subsection{Windows Server --- Fase 2 (Desktop Experience)}

Instalación de la versión gráfica (Desktop Experience). Evidencias:

\subsubsection{Instalación}

\begin{figure}[H]
  \centering
  \includegraphics[width=0.9\linewidth]{windows-gui-01-instalando.png}
  \caption{Instalación de Windows Server 2025 (edición Desktop Experience) en curso: copia de archivos.}
  \label{fig:wingui-01}
\end{figure}

OOBE (Out-of-Box Experience) completado: cuenta \texttt{Administrator} con contraseña \texttt{Server2025!}.

Red: estática aplicada con \texttt{netsh} (10.2.78.67/16, GW y DNS 10.2.65.1). Pruebas OK (0\,\% de pérdida): gateway, 8.8.8.8, \texttt{www.google.com} e inter-VM con Slackware (.64).

\subsubsection{Primer arranque (escritorio y Server Manager)}

\begin{figure}[H]
  \centering
  \includegraphics[width=0.9\linewidth]{windows-gui-02-escritorio-server-manager.png}
  \caption{Escritorio de Windows Server 2025 con Server Manager abierto (resumen del servidor local y aviso de Azure Arc).}
  \label{fig:wingui-02}
\end{figure}

\textbf{Usuarios y grupos (evidencia: \texttt{lab-evidencias/usuarios-windows/\allowbreak usuarios-creados.txt})}

Se crearon los 4 usuarios con PowerShell (\texttt{New-LocalUser}), todos habilitados y con la descripción significativa pedida por el lab; los usuarios 1 y 2 se asignaron al grupo local \texttt{Accounting} y los usuarios 3 y 4 a \texttt{IT} (\texttt{New-LocalGroup} + \texttt{Add-LocalGroupMember}):

\begin{lstlisting}[language={}]
New-LocalUser -Name claudia -Description "Usuario con el nombre de la profesora"
New-LocalUser -Name john    -Description "Usuario con el nombre del profesor"
New-LocalUser -Name fabian  -Description "Usuario con el nombre del profesor"
New-LocalUser -Name diego   -Description "Usuario con el nombre del profesor"

New-LocalGroup -Name Accounting
New-LocalGroup -Name IT

Add-LocalGroupMember -Group Accounting -Member claudia, john
Add-LocalGroupMember -Group IT -Member fabian, diego
\end{lstlisting}

Resultado verificado (05/08/2026): \texttt{claudia}, \texttt{john}, \texttt{fabian} y \texttt{diego} aparecen en la lista de usuarios locales con \texttt{Enabled = True}; el grupo \texttt{Accounting} contiene a \texttt{WIN-3VMPCF0KCDN\textbackslash{}claudia} y \texttt{WIN-3VMPCF0KCDN\textbackslash{}john}, y el grupo \texttt{IT} a \texttt{WIN-3VMPCF0KCDN\textbackslash{}fabian} y \texttt{WIN-3VMPCF0KCDN\textbackslash{}diego}.

\textbf{Permisos con ACL e \texttt{icacls} (evidencia: \texttt{lab-evidencias/usuarios-windows/\allowbreak permisos-acl.txt})}

Sobre la carpeta \texttt{C:\textbackslash{}LabPermisos} se asignaron niveles de permiso con \texttt{icacls} (con \texttt{/grant} y \texttt{/deny}) y se verificó la ACL final:

\begin{table}[H]
  \centering
  \begin{tabular}{@{}lll@{}}
    \toprule
    \textbf{Usuario} & \textbf{Nivel} & \textbf{Significado} \\
    \midrule
    \texttt{claudia} & \texttt{(OI)(CI)(F)} & Control total (Full) \\
    \texttt{john}    & \texttt{(OI)(CI)(R)} & Solo lectura (Read) \\
    \texttt{fabian}  & \texttt{(OI)(CI)(M)} & Modificar (Modify) \\
    \texttt{diego}   & \texttt{(OI)(CI)(N)} & Deny (denegar) \\
    \bottomrule
  \end{tabular}
  \caption{Niveles de permiso asignados sobre \texttt{C:\textbackslash{}LabPermisos}.}
  \label{tab:win-acl}
\end{table}

\texttt{(OI)} = heredado por objetos (archivos) y \texttt{(CI)} = heredado por contenedores (subcarpetas). El propietario de la carpeta es \texttt{BUILTIN\textbackslash{}Administrators}; el archivo de prueba \texttt{prueba.txt} hereda la misma matriz (marcas \texttt{(I)}): claudia F, fabian M, john R y diego N. En Windows la negación (Deny) tiene prioridad sobre los permisos concedidos, por lo que \texttt{diego} queda excluido del acceso.

\textbf{Visor de eventos --- eventos 4624 y 4625 (evidencia: \texttt{lab-evidencias/usuarios-windows/\allowbreak eventos-seguridad.txt})}

\begin{itemize}
  \item \textbf{4625 (logon fallido):} un intento con contraseña incorrecta de \texttt{claudia} el 05/08/2026 19:58:49 y dos intentos sobre \texttt{Administrator} el 03/08/2026, con estado \texttt{\%\%2313} (credenciales inválidas).
  \item \textbf{4624 (logon exitoso):} \texttt{claudia} el 05/08/2026 19:59:02, \texttt{Administrator} a las 19:58:25 y la cuenta \texttt{SYSTEM} a las 19:59:07.
\end{itemize}

\textbf{Windows Registry (evidencia: \texttt{lab-evidencias/usuarios-windows/registry.txt} y Figura~\ref{fig:win-registry})}

\begin{itemize}
  \item Qué es: base de datos jerárquica donde Windows y las aplicaciones guardan su configuración, organizada en colmenas (\emph{hives}). La colmena principal es HKLM (\texttt{HKEY\_LOCAL\_MACHINE}), con las claves \texttt{BCD00000000}, \texttt{HARDWARE}, \texttt{SAM}, \texttt{SECURITY}, \texttt{SOFTWARE} y \texttt{SYSTEM}.
  \item Ejemplos verificados: el hostname está en \texttt{HKLM\allowbreak\textbackslash{}SYSTEM\allowbreak\textbackslash{}CurrentControlSet\allowbreak\textbackslash{}Control\allowbreak\textbackslash{}ComputerName\allowbreak\textbackslash{}ComputerName} $\rightarrow$ \texttt{WIN-3VMPCF0KCDN}; la ruta de instalación en \texttt{HKLM\allowbreak\textbackslash{}SOFTWARE\allowbreak\textbackslash{}Microsoft\allowbreak\textbackslash{}Windows\allowbreak\textbackslash{}CurrentVersion\allowbreak\textbackslash{}ProgramFilesDir} $\rightarrow$ \texttt{C:\textbackslash{}Program Files}.
  \item Edición con \texttt{regedit} (captura: Figura~\ref{fig:win-registry}): la herramienta gráfica muestra el árbol de claves y valores; se editan con doble clic (también se puede en línea de comandos con \texttt{reg add} / \texttt{reg query}).
\end{itemize}

\begin{figure}[H]
  \centering
  \includegraphics[width=0.85\linewidth]{usuarios-windows/registry-editor.png}
  \caption{Edición del registro con \texttt{regedit} en Windows Server 2025: árbol de claves de HKLM.}
  \label{fig:win-registry}
\end{figure}

\subsection{Android}

Creación de la VM e instalación de \textbf{android-x86 9.0-r2} (kernel 4.9). Evidencias:

\subsubsection{Instalación}

\begin{figure}[H]
  \centering
  \includegraphics[width=0.9\linewidth]{android-01-menu-instalacion.png}
  \caption{Menú de arranque de Android-x86 9.0-r2: \texttt{Live CD}, instalación en disco y opciones avanzadas.}
  \label{fig:and-01}
\end{figure}

\begin{figure}[H]
  \centering
  \includegraphics[width=0.9\linewidth]{android-02-instalador-elegir-particion.png}
  \caption{Selección de la partición de destino para instalar Android-x86 (disco con 8 GB de espacio libre).}
  \label{fig:and-02}
\end{figure}

\begin{figure}[H]
  \centering
  \includegraphics[width=0.9\linewidth]{android-03-confirm-gpt.png}
  \caption{Confirmación del esquema de particionado: \texttt{Do you want to use GPT?}}
  \label{fig:and-03}
\end{figure}

\begin{figure}[H]
  \centering
  \includegraphics[width=0.9\linewidth]{android-04-cfdisk.png}
  \caption{Particionado del disco con \texttt{cfdisk}: disco \texttt{/dev/sdd} de 17,1 GB.}
  \label{fig:and-04}
\end{figure}

\begin{figure}[H]
  \centering
  \includegraphics[width=0.9\linewidth]{android-05-particion-creada.png}
  \caption{Partición creada y visible en la pantalla de selección para la instalación.}
  \label{fig:and-05}
\end{figure}

Selección de la partición para la instalación (incluido el modo automático):

\begin{figure}[H]
  \centering
  \includegraphics[width=0.9\linewidth]{android-06-choose-partition-v2.png}
  \caption{Selección de la partición de destino para la instalación (segunda vista de la misma pantalla de elección).}
  \label{fig:and-06}
\end{figure}

\begin{figure}[H]
  \centering
  \includegraphics[width=0.9\linewidth]{android-07-choose-partition-auto.png}
  \caption{Selección de la partición en modo automático (confirmación del destino de instalación).}
  \label{fig:and-07}
\end{figure}

\begin{figure}[H]
  \centering
  \includegraphics[width=0.9\linewidth]{android-08-instalado-exitoso.png}
  \caption{Instalación completada con éxito: Android-x86 9.0-r2 instalado en el disco.}
  \label{fig:and-08}
\end{figure}

\begin{figure}[H]
  \centering
  \includegraphics[width=0.9\linewidth]{android-09-home-app-select.png}
  \caption{Selección de la aplicación de inicio (Home): \texttt{Quickstep} o \texttt{Taskbar}.}
  \label{fig:and-09}
\end{figure}

\subsubsection{Primer arranque y escritorio}

Escritorio de Android funcionando:

\begin{figure}[H]
  \centering
  \includegraphics[width=0.9\linewidth]{android-10-escritorio.png}
  \caption{Escritorio de Android funcionando (1).}
  \label{fig:and-10}
\end{figure}

\begin{figure}[H]
  \centering
  \includegraphics[width=0.9\linewidth]{android-11-escritorio-final.png}
  \caption{Escritorio de Android funcionando (vista final).}
  \label{fig:and-11}
\end{figure}

Datos de la VM: 2 CPU, 2 GB de RAM, disco VDI de 16 GB, adaptador de red NAT y GRUB instalado en el disco (arranque directo desde disco). El arranque presentó un problema de visualización que se resolvió habilitando 128 MB de VRAM y la aceleración 3D en la configuración de la VM.

\textbf{Red (detalle en el párrafo ``Pruebas de conectividad en la red de la universidad'' de la sección 3.2.4)}

IP estática \texttt{10.2.78.68/16} sobre la interfaz ethernet \texttt{wifi\_eth}, con \emph{policy routing} en la tabla 97 (reglas de \texttt{ip rule}/\texttt{ip route} para salir por el gateway 10.2.65.1 y alcanzar la red 10.2.0.0/16). Pings OK al gateway, a 8.8.8.8 e inter-VM con Slackware (.64). Limitación documentada: la resolución DNS por nombre no funciona desde la consola de emergencia (Android sin \texttt{netd}). Evidencia: \texttt{lab-evidencias/red/android-pings-uni2.png}.

\subsection{Conocimiento de línea de comandos}

\begin{itemize}
  \item ¿Qué es el shell? Intérprete de comandos que ejecuta programas y da acceso a los servicios del sistema (login, tuberías, variables, scripts). En Unix es a la vez prompt interactivo y lenguaje de scripting; en Windows conviven \texttt{cmd} (interpretador heredado de MS-DOS) y PowerShell (orientado a objetos).
  \item Shells soportados, verificados con la lista de shells válidos del sistema (evidencia \texttt{lab-evidencias/shell/}):
    \begin{itemize}
      \item \textbf{Slackware 15.0} (\texttt{slackware-comandos.txt}): bash, ksh, csh, tcsh, zsh, dash y ash (más \texttt{sh}); los ejecutables están en \texttt{/bin} y \texttt{/usr/bin}.
      \item \textbf{Solaris 11.4} (\texttt{solaris-comandos.txt}): bash, ksh, csh, zsh y sh (\texttt{/usr/bin/bash}, \texttt{/usr/bin/ksh}, \texttt{/usr/bin/csh}, \texttt{/usr/bin/zsh}, \texttt{/usr/bin/sh}).
      \item \textbf{Windows Server 2025} (\texttt{windows-core-comandos.txt}): \texttt{cmd.exe} (Windows 10.0.26100) y PowerShell 5.1 (5.1.26100.1591).
    \end{itemize}
  \item Diferencias: en Unix cada shell es un programa distinto con su propia sintaxis; Solaris usa históricamente sh/ksh/csh y Slackware incluye la familia completa. En Windows, \texttt{cmd} manipula texto plano y PowerShell trabaja con objetos; los homólogos de los comandos Unix (\texttt{dir}, \texttt{copy}, \texttt{findstr}, \texttt{where}) son menos potentes para procesar texto.
\end{itemize}

Comandos (tabla comparativa Linux/Unix vs.\ Windows):

\begin{table}[H]
  \centering
  \begin{tabular}{@{}lll@{}}
    \toprule
    \textbf{Tarea} & \textbf{Linux/Unix} & \textbf{Windows} \\
    \midrule
    Cambiar de directorio & \texttt{cd} & \texttt{cd} \\
    Listar archivos & \texttt{ls -la} & \texttt{dir} \\
    Copiar/mover archivo & \texttt{cp}, \texttt{mv} & \texttt{copy}, \texttt{move} \\
    Ver contenido sin editar & \texttt{cat} & \texttt{type} \\
    Editar un archivo & \texttt{vi}/\texttt{nano} & \texttt{notepad} \\
    Primeras/últimas líneas & \texttt{head}, \texttt{tail} & \texttt{Get-Content -First/-Last} (PowerShell) \\
    Buscar palabra en archivo & \texttt{grep} & \texttt{findstr} / \texttt{Select-String} \\
    Localizar un archivo & \texttt{find}, \texttt{locate} & \texttt{where} / \texttt{Get-ChildItem -Recurse} \\
    \bottomrule
  \end{tabular}
  \caption{Comparativa de comandos Linux/Unix vs.\ Windows.}
  \label{tab:shells}
\end{table}

Ejemplos de uso de los 8 comandos con salidas reales abreviadas (evidencia \texttt{lab-evidencias/shell/}):

\textbf{Slackware 15.0} (prompt \texttt{\$}):

\begin{verbatim}
$ cd /tmp
/tmp

$ ls -la /usuarios
claudia (accounting), john, fabian, diego

$ mv archivo.txt renombrado.txt
-rwxr----- 1 root root /tmp/renombrado.txt

$ cat archivo.txt
slackware

$ echo "linea-anadida" >> archivo.txt
linea-anadida

$ tail /etc/passwd
root:x:0:0:...
diego:x:1004:...

$ grep claudia /etc/passwd
claudia:x:1001:1002:...

$ find /usuarios -name .screenrc
4 archivos
\end{verbatim}

\textbf{Solaris 11.4} (prompt \texttt{\$}):

\begin{verbatim}
$ cd /tmp
/tmp

$ ls -la /usuarios
claudia (Accounting), john, fabian, diego

$ mv archivo.txt renombrado.txt
-rwxr----- 1 root root /tmp/renombrado.txt

$ cat archivo.txt
solaris

$ echo "linea-anadida" >> archivo.txt
linea-anadida

$ tail /etc/passwd
root:x:0:0:Super-User
diego:x:103:...

$ grep claudia /etc/passwd
claudia:x:100:100:...

$ find /usuarios -name .profile
.profile de cada usuario
\end{verbatim}

\textbf{Windows Server 2025} (prompt \texttt{C:\textbackslash{}>}):

\begin{verbatim}
C:\> cd \
C:\

C:\> dir C:\LabPermisos
prueba.txt  20 bytes

C:\> copy archivo.txt copia.txt
1 file(s) copied.

C:\> move archivo.txt renombrado.txt
renombrado.txt

C:\> type archivo.txt
archivo-de-prueba

C:\> echo linea-anadida >> archivo.txt
linea-anadida

C:\> Get-Content -Last hosts
localhost

C:\> findstr localhost hosts
localhost

C:\> where notepad
C:\Windows\System32\notepad.exe
\end{verbatim}

Nota: los ejemplos se ejecutaron y verificaron en los tres sistemas.

\clearpage
\section{Uso y aplicaciones}

\begin{itemize}
  \item \textbf{Administración de servidores Unix y Windows:} lo aprendido permite administrar servidores sin interfaz gráfica: configuración de red estática, usuarios y grupos, servicios y línea de comandos. Slackware y Windows Server Core se administran por consola (SSH y SConfig), el modelo que se encuentra en servidores de archivos, DNS, correo y bases de datos reales.
  \item \textbf{Monitoreo con logs:} syslog (Slackware y Solaris) y el Visor de eventos (Windows) permiten detectar fallos de hardware y red, errores de servicios (sendmail, sshd) e intentos de acceso no autorizados (eventos 4624/4625, log de sudo), base de la auditoría de seguridad.
  \item \textbf{Gestión de usuarios y permisos:} la creación de usuarios y grupos con home en \texttt{/usuarios} y los permisos por grupo (\texttt{chmod}/ACL/\texttt{icacls}) aplican el principio de menor privilegio: en una empresa cada departamento (Accounting/IT) accede solo a sus recursos.
  \item \textbf{Virtualización en infraestructura empresarial y cloud:} un solo servidor físico consolidado en varias VM (menos hardware, menos energía), entornos de prueba aislados, instantáneas y migración entre hipervisores (OVF/OVA), y la base del modelo IaaS de las nubes públicas: el hipervisor es el componente sobre el que se construye la elasticidad y la multitenencia; los contenedores lo complementan para los microservicios.
\end{itemize}

\section{Conclusiones}

\textit{Borrador provisional; se ajustará al cerrar el laboratorio.}

\begin{itemize}
  \item Se instalaron y dejaron operativos cinco sistemas sobre un único equipo: Slackware 15.0, Solaris 11.4, Windows Server 2025 (Core y Desktop Experience) y Android-x86 9.0-r2, cada uno con su instalador, su sistema de archivos y su modelo de administración.
  \item La virtualización demostró su valor práctico: consolidación de plataformas heterogéneas, recuperación ante errores (reinstalaciones y snapshots) y portabilidad entre hipervisores (formatos OVF).
  \item La red estática del laboratorio (sin DHCP) obligó a planificar las direcciones y a configurar cada equipo a mano; la conectividad quedó verificada con 0\,\% de pérdida en todas las pruebas de Slackware, Windows Core, Windows GUI y Android.
  \item Se documentaron particularidades reales de cada plataforma: Windows bloquea el ICMP entrante por firewall, Android requiere policy routing y no resuelve DNS por nombre desde la consola de emergencia, y Slackware impone limitaciones en los nombres de grupo.
  \item Los conceptos de permisos, logs y shells funcionan de forma equivalente entre Slackware y Solaris, con diferencias de implementación (ACL/ZFS, SMF, IPS).
  \item Pendiente: verificar RDP en Windows Server Core y ajustar estas conclusiones con los resultados finales.
\end{itemize}

\section{Bibliografía}

\begin{itemize}
  \item Oracle. \emph{Oracle Solaris 11.4 Documentation}. \url{https://docs.oracle.com/en/solaris/solaris/11.4/}
  \item Slackware Linux Project. \emph{Slackware Linux 15.0 Documentation}. \url{https://www.slackware.com/documentation/}
  \item Microsoft. \emph{Evaluation Center --- Windows Server 2025}. \url{https://www.microsoft.com/en-us/evalcenter/}
  \item Android-x86 Project. \emph{Android-x86 --- Release 9.0-r2}. \url{https://www.android-x86.org/}
  \item Oracle. \emph{Oracle VM VirtualBox Documentation}. \url{https://www.virtualbox.org/wiki/Docs}
  \item VMware (Broadcom). \emph{VMware Docs --- máquinas virtuales y conversión OVF}. \url{https://docs.vmware.com/}
  \item The Linux Foundation. \emph{Filesystem Hierarchy Standard 3.0}. \url{https://refspecs.linuxfoundation.org/FHS_3.0/fhs-3.0.pdf}
\end{itemize}

\section*{Anexo A --- Asignación de IPs (10.2.78.n)}
\addcontentsline{toc}{section}{Anexo A --- Asignación de IPs (10.2.78.n)}

La red del laboratorio es 10.2.0.0/16, con gateway y DNS 10.2.65.1 y dominio \texttt{is.escuelaing.edu.co}. \textbf{NO tiene DHCP}: todas las direcciones se configuraron estáticamente en cada equipo.

\begin{table}[H]
  \centering
  \begin{tabularx}{\linewidth}{@{}lXp{5.5cm}@{}}
    \toprule
    \textbf{Máquina} & \textbf{IP} & \textbf{Notas} \\
    \midrule
    Host (portátil físico) & \texttt{10.2.78.69} & Estática (configurada con \texttt{netsh} sobre el adaptador Ethernet), GW y DNS 10.2.65.1 \\
    Slackware & \texttt{10.2.78.64} & Estática (\texttt{rc.inet1.conf}), probada OK \\
    Solaris & \texttt{10.2.78.65} & Estática (\texttt{ipadm}), GW y DNS 10.2.65.1 \\
    Windows Server Core & \texttt{10.2.78.66} & Estática (\texttt{netsh}), probada OK; ICMP habilitado en el firewall \\
    Windows Server GUI & \texttt{10.2.78.67} & Estática (\texttt{netsh}), probada OK \\
    Android & \texttt{10.2.78.68} & Estática en \texttt{wifi\_eth} con policy routing (tabla 97); DNS por nombre no resuelve \\
    PC del laboratorio (referencia) & \texttt{10.2.254.6} & Equipo físico del aula, gateway 10.2.65.1 \\
    \bottomrule
  \end{tabularx}
  \caption{Asignación de direcciones IP (10.2.78.n).}
  \label{tab:ips}
\end{table}

\section*{Anexo B --- Guion del video (punto 1)}
\addcontentsline{toc}{section}{Anexo B --- Guion del video (punto 1)}

Video del grupo: \url{https://youtu.be/6HnwGg-_QV0}

Guion para el video de hasta 5 minutos sobre virtualización. Estilo presentación hablada; cada bloque indica el tiempo aproximado.

\textbf{(0:00--0:30) Apertura}

Hola, y bienvenidos. En este video vamos a explicar qué es la virtualización, cómo funcionan los hipervisores, por qué son la base del cómputo en la nube y en qué se diferencian las máquinas virtuales de los contenedores. La idea central es sencilla: un solo computador físico puede ejecutar varios sistemas operativos a la vez, cada uno aislado de los demás, como si fueran computadores independientes.

\textbf{(0:30--1:15) Hipervisores: qué son y cómo se clasifican}

El software que hace posible esto es el hipervisor, también llamado monitor de máquina virtual. Es una capa que se ubica entre el hardware y los sistemas operativos invitados: reparte la CPU, la memoria, el disco y la red entre las máquinas, y garantiza que una máquina no interfiera con otra. Los hipervisores se clasifican en dos tipos. El Tipo 1, o bare-metal, se instala directamente sobre el hardware, sin un sistema operativo debajo: ejemplos son VMware ESXi, Microsoft Hyper-V, Xen y KVM, y es el modelo que usan los centros de datos. El Tipo 2, o hosted, corre como una aplicación sobre un sistema operativo anfitrión: VirtualBox y VMware Workstation, que es lo que usamos en el laboratorio para instalar Slackware, Solaris, Windows y Android.

\textbf{(1:15--2:00) Características y arquitectura}

Las características clave de un hipervisor son el aislamiento, la partición de recursos y la encapsulación. Aislamiento: si una máquina virtual falla o se infecta, las demás siguen funcionando. Partición: el hipervisor divide los recursos físicos entre todas las máquinas. Encapsulación: una máquina virtual completa es un conjunto de archivos, así que se puede copiar, respaldar o migrar como cualquier otro archivo. En cuanto a la arquitectura, cada máquina virtual cree tener su propio hardware: una CPU virtual, memoria virtual, un disco y una tarjeta de red virtuales. El hipervisor intercepta las instrucciones privilegiadas del sistema invitado y las ejecuta sobre el hardware real, apoyándose en las extensiones de virtualización de los procesadores actuales.

\textbf{(2:00--2:45) Cloud computing y su relación con los hipervisores}

Ahora, ¿qué es el cómputo en la nube? Es la entrega de recursos de cómputo, almacenamiento y software a través de internet, bajo demanda y con pago por uso. Hay tres modelos de servicio: IaaS, que ofrece máquinas virtuales y redes; PaaS, que ofrece la plataforma para desplegar aplicaciones; y SaaS, que ofrece el software completo. ¿Se usan hipervisores en la nube? Sí, definitivamente: cuando en AWS o Azure pides una máquina virtual, es un hipervisor el que la crea en segundos sobre un servidor físico compartido con otros clientes. La multitenencia y la elasticidad, las dos propiedades que definen a la nube, son posibles gracias a la virtualización.

\textbf{(2:45--3:15) Costos: servidor físico vs.\ nube}

En costos hay un tradeoff. Un servidor físico exige una inversión inicial alta en hardware y costos fijos de energía, enfriamiento y mantenimiento, y suele quedar sobredimensionado para los picos de carga. La nube elimina esa inversión: se paga por uso, sin mantenimiento y con elasticidad real. Pero el costo es continuo: si el servidor se usa al máximo todo el tiempo, a largo plazo la nube puede salir más cara que el hardware propio. La decisión depende del perfil de carga y del nivel de control que se necesite.

\textbf{(3:15--4:15) Contenedores vs.\ máquinas virtuales}

Finalmente, hablemos de los contenedores. Una máquina virtual virtualiza el hardware: cada VM carga su propio sistema operativo completo. Un contenedor, en cambio, virtualiza el sistema operativo: todos los contenedores comparten el kernel del anfitrión, y el aislamiento se logra con namespaces y cgroups, las herramientas del kernel de Linux. Por eso un contenedor pesa megabytes y arranca en segundos, mientras una máquina virtual pesa gigabytes y arranca en minutos; por eso en un servidor caben decenas de máquinas virtuales pero cientos o miles de contenedores. La contrapartida es el aislamiento: la frontera de una máquina virtual es mucho más fuerte que la de un contenedor. En la práctica no compiten: las máquinas virtuales aíslan plataformas completas y los contenedores despliegan microservicios sobre ellas.

\textbf{(4:15--4:45) Cierre}

En resumen: el hipervisor es el software que multiplica un servidor físico en varias máquinas virtuales; la nube pública se construye sobre hipervisores; los contenedores comparten el kernel y son la opción ligera para las aplicaciones. En el laboratorio aplicamos estos conceptos instalando y administrando varios sistemas sobre una sola computadora, y verificando su conectividad en una red real. Gracias por ver el video.

\section*{Anexo C --- Conversión a LaTeX}
\addcontentsline{toc}{section}{Anexo C --- Conversión a LaTeX}

\begin{itemize}
  \item Mapeo 1:1 de secciones a \texttt{\textbackslash section} / \texttt{\textbackslash subsection}.
  \item Tablas $\rightarrow$ \texttt{tabular} / \texttt{booktabs}; capturas $\rightarrow$ \texttt{\textbackslash includegraphics} con \texttt{lab-evidencias/} como ruta.
\end{itemize}

\end{document}
